% !PRINTED 30
\weberpage{30}
Um aber nachzuweisen, dass $R$ eine commutative Gruppe ist, müssen wir zwei Hülfsbetrachtungen voranschicken.

1) Es ist zu beweisen, dass irgend zwei der Gruppen $Q_\nu, Q'_\nu, Q''_\nu \ldots$ ausser der Einheit keinen gemeinschaftlichen Theiler haben.

Wir brauchen hierzu nur nachzuweisen, dass $Q_\nu, Q'_\nu$ diese Eigenschaft haben. Denn haben
\[ Q'_\nu = \gamma'^{-1} Q_\nu \gamma', \quad Q''_\nu = \gamma''^{-1} Q_\nu \gamma'' \]
den grössten gemeinschaftlichen Theiler $T$, so haben $Q_\nu$ und $Q'''_\nu = \gamma' Q''_\nu \gamma'^{-1}$ den grössten gemeinschaftlichen Theiler $\gamma' T \gamma'^{-1}$.

Nehmen wir also an, es sei $T$ der grösste gemeinschaftliche Theiler von $Q_\nu$ und $Q'_\nu$, so ist $T$ gewiss eine commutative Gruppe, weil $Q_\nu, Q'_\nu$ solche sind. Es ist aber $T$ auch Normaltheiler von $P_\nu$. Denn $Q_\nu$ und $Q'_\nu$ sind solche Theiler; und wenn also $\tau$ ein Element aus $T$ und $\pi$ ein Element aus $P_\nu$ ist, so ist $\pi^{-1}\tau\pi$ sowohl in $Q_\nu$ als in $Q'_\nu$, also auch in $T$ enthalten. Folglich ist $\pi^{-1} T \pi = T$.

Nun haben wir aber angenommen, dass $Q_\nu$ ein commutativer Normaltheiler von $P_\nu$ über der Einheit von möglichst niedrigem Grade sei; ferner waren $Q_\nu$ und $Q'_\nu$ von einander verschieden. Also kann $T$ nur die Einheitsgruppe selbst sein, w. z. b. w.

2) Um die Gruppe $R$ zu bilden, können wir so verfahren, dass wir die Elemente von $Q_\nu, Q'_\nu, Q''_\nu \ldots$ in beliebiger Reihenfolge und Wiederholung so lange mit einander componiren, als noch neue Elemente entstehen. Denn alle so gebildeten Zusammensetzungen bilden eine Gruppe $C$, die in $R$ enthalten ist, weil $R$ die einzelnen Elemente von $Q_\nu, Q'_\nu, Q''_\nu \ldots$ enthält. Andererseits sind auch die Gruppen $Q_\nu, Q'_\nu, Q''_\nu \ldots$ in $C$ enthalten, und folglich ist $R$ in $C$ enthalten und $C$ ist mit $R$ identisch.

Um also endlich zu beweisen, dass $R$ eine commutative Gruppe ist, bleibt uns nur zu zeigen, dass, wenn $\alpha, \beta$ irgend zwei Elemente aus $Q_\nu, Q'_\nu, Q''_\nu \ldots$ sind, die Vertauschbarkeit
\weq{10}{\alpha\beta = \beta\alpha}
besteht.

Wenn nun $\alpha, \beta$ beide in dieselbe Gruppe $Q_\nu$ gehören, so ist (10) Folge unserer Annahme, dass $Q_\nu$ commutativ sei; und ebenso ist es, wenn beide in einer der anderen Gruppen $Q'_\nu, Q''_\nu \ldots$ vorkommen.
